Bei einem adiabatischen Prozess gilt für ein zweiatomiges Gas die Beziehung
\[
 p\cdot V^{\frac{7}{5}} = \text{konstant}
\]
\begin{abcliste}
\abc Bestimmen Sie mit dem Taschenrechner anhand einer graphischen Darstellung die Arbeit, die bei der adiabatischen Expansion von Sauerstoff vom Anfangszustand (\paO, \VaO) aufs doppelte Volumen verrichtet wird.
\abc Berechnen Sie die Arbeit näherungsweise, indem Sie den Anfangs- und den Endzustand des Prozesses im $p(V)$-Diagramm durch ein Geradenstück verbinden.
\end{abcliste}
